\section{引言}
\label{ch:intro}

\subsection{Agent Evolution的定义与范畴}

Agent Evolution是指以LLM或多模态基础模型为核心推理器，以环境反馈、任务奖励、人类偏好、程序化测试为改进信号，通过可重复的``生成—评估—选择—更新''循环，使智能体系统产生可验证变化的一类方法。

% ====================================================================
% 图1.1: Agent Evolution 六大方法分类体系
% ====================================================================
\begin{figure}[htbp]
\centering
\begin{tikzpicture}[
    root/.style={rectangle, rounded corners=8pt, draw=blue!60!black, fill=blue!15,
        text width=4.5cm, minimum height=1cm, align=center, font=\bfseries\large},
    cat/.style={rectangle, rounded corners=6pt, draw=#1!70!black, fill=#1!15,
        text width=3.8cm, minimum height=2.8cm, align=center, font=\small},
    example/.style={font=\scriptsize, text width=3.4cm, align=left},
    edge/.style={-{Stealth[length=3mm]}, thick, color=gray!60},
]
    % Root
    \node[root] (ae) {Agent Evolution\\自演化方法体系};

    % Categories - 2 rows of 3
    \node[cat=cat1] (c1) at (-6,-4) {};
    \node[cat=cat2] (c2) at (-2,-4) {};
    \node[cat=cat3] (c3) at (2,-4) {};
    \node[cat=cat4] (c4) at (-6,-8.5) {};
    \node[cat=cat5] (c5) at (-2,-8.5) {};
    \node[cat=cat6] (c6) at (2,-8.5) {};

    % Labels inside categories
    \node[font=\bfseries\small, color=cat1!70!black] at (-6,-3) {基于奖励的进化};
    \node[example, color=cat1!50!black] at (-6,-3.7) {STaR, Self-Rewarding,\\Meta-Rewarding, RISE,\\Agent-R, RAGEN};

    \node[font=\bfseries\small, color=cat2!70!black] at (-2,-3) {自博弈进化};
    \node[example, color=cat2!50!black] at (-2,-3.7) {Self-Play FT, SPIRAL,\\Absolute Zero,\\Multi-Agent Debate};

    \node[font=\bfseries\small, color=cat3!70!black] at (2,-3) {提示词进化};
    \node[example, color=cat3!50!black] at (2,-3.7) {Self-Refine, Reflexion,\\ExpeL, ACE, EvolveR,\\符号学习};

    \node[font=\bfseries\small, color=cat4!70!black] at (-6,-7.5) {架构搜索};
    \node[example, color=cat4!50!black] at (-6,-8.2) {ADAS, G\"{o}del Agent,\\SICA, DGM,\\AlphaEvolve, EvoMAC};

    \node[font=\bfseries\small, color=cat5!70!black] at (-2,-7.5) {基于记忆的进化};
    \node[example, color=cat5!50!black] at (-2,-8.2) {Voyager, Generative\\Agents, ReasoningBank,\\Memory-R1, AriadneMem};

    \node[font=\bfseries\small, color=cat6!70!black] at (2,-7.5) {混合方法};
    \node[example, color=cat6!50!black] at (2,-8.2) {WebEvolver,\\开放式进化闭环,\\多机制组合系统};

    % Edges
    \draw[edge] (ae.south) -- ++(0,-0.5) -| (c1.north);
    \draw[edge] (ae.south) -- ++(0,-0.5) -| (c2.north);
    \draw[edge] (ae.south) -- ++(0,-0.5) -| (c3.north);
    \draw[edge] (ae.south) -- ++(0,-0.5) -| (c4.north);
    \draw[edge] (ae.south) -- ++(0,-0.5) -| (c5.north);
    \draw[edge] (ae.south) -- ++(0,-0.5) -| (c6.north);

\end{tikzpicture}
\caption{Agent Evolution六大方法分类体系。每个类别依据``选择压力''和``可变对象''划分，覆盖从奖励驱动到架构搜索的完整光谱。}
\label{fig:taxonomy}
\end{figure}

% ====================================================================
% 图1.2: 技术演进时间线
% ====================================================================
\begin{figure}[htbp]
\centering
\begin{tikzpicture}[
    year/.style={font=\bfseries\large, color=blue!70!black},
    paper/.style={rectangle, rounded corners=3pt, draw=#1!60!black, fill=#1!10,
        font=\scriptsize, text width=2.8cm, minimum height=0.6cm, align=center, inner sep=2pt},
    timeline/.style={very thick, -{Stealth[length=4mm]}, color=gray!70},
]
    % Timeline
    \draw[timeline] (0,0) -- (16,0);
    \foreach \x/\y in {1/2022, 4.5/2023, 8/2024, 11.5/2025, 15/2026} {
        \draw[thick, color=gray!50] (\x,-0.2) -- (\x,0.2);
        \node[year, below] at (\x,-0.4) {\y};
    }

    % 2022
    \node[paper=cat1] (star) at (1,1.5) {STaR\\推理链自举};
    \draw[-{Stealth}, color=cat1!50] (star.south) -- (1,0.2);

    % 2023
    \node[paper=cat3] (ref) at (3.5,1.5) {Reflexion\\语言反思};
    \node[paper=cat3] (sr) at (3.5,3) {Self-Refine\\迭代精炼};
    \node[paper=cat5] (voy) at (5.5,1.5) {Voyager\\技能库};
    \draw[-{Stealth}, color=cat3!50] (ref.south) -- (3.5,0.2);
    \draw[-{Stealth}, color=cat3!50] (sr.south) -- (3.5,0.2);
    \draw[-{Stealth}, color=cat5!50] (voy.south) -- (5.5,0.2);

    % 2024
    \node[paper=cat1] (selfr) at (7,1.5) {Self-Rewarding\\自奖励LM};
    \node[paper=cat4] (adas) at (7,3) {ADAS\\架构搜索};
    \node[paper=cat4] (godel) at (9,1.5) {G\"{o}del Agent\\自修改};
    \node[paper=cat3] (evo) at (9,3) {EvoMAC\\多Agent文本BP};
    \draw[-{Stealth}, color=cat1!50] (selfr.south) -- (7,0.2);
    \draw[-{Stealth}, color=cat4!50] (adas.south) -- (7,0.2);
    \draw[-{Stealth}, color=cat4!50] (godel.south) -- (9,0.2);
    \draw[-{Stealth}, color=cat3!50] (evo.south) -- (9,0.2);

    % 2025
    \node[paper=cat4] (dgm) at (11,1.5) {DGM\\Darwin G\"{o}del};
    \node[paper=cat4] (ae) at (11,3) {AlphaEvolve\\算法发现};
    \node[paper=cat2] (az) at (13,1.5) {Absolute Zero\\零数据自博弈};
    \node[paper=cat3] (ace) at (13,3) {ACE\\上下文工程};
    \draw[-{Stealth}, color=cat4!50] (dgm.south) -- (11,0.2);
    \draw[-{Stealth}, color=cat4!50] (ae.south) -- (11,0.2);
    \draw[-{Stealth}, color=cat2!50] (az.south) -- (13,0.2);
    \draw[-{Stealth}, color=cat3!50] (ace.south) -- (13,0.2);

    % 2026
    \node[paper=cat5] (ariadne) at (15,1.5) {AriadneMem\\记忆架构};
    \draw[-{Stealth}, color=cat5!50] (ariadne.south) -- (15,0.2);

\end{tikzpicture}
\caption{Agent Evolution关键技术演进时间线（2022--2026）。颜色对应六大方法分类：蓝色（奖励）、橙色（自博弈）、绿色（提示词）、紫色（架构搜索）、红色（记忆）、青色（混合）。}
\label{fig:timeline}
\end{figure}

\subsection{研究背景与动机}

Agent Evolution兴起的五个核心背景：（1）基础模型能力增长与静态调用范式的矛盾；（2）人工监督成本上升；（3）真实场景要求长期适应；（4）安全与可控性问题；（5）进化计算与LLM的互补性。

\subsection{综述结构与阅读指南}

本综述按``理论—方法—系统—评估—实践—痛点—前沿''组织。读者可从第3章（方法分类）作为快速入口，结合图\ref{fig:taxonomy}的六类方法和图\ref{fig:timeline}的时间线定位感兴趣的工作。
